\section{Discussion and open problems}\label{sec:discussion}

\paragraph{Summary.}
A loss ratio carries its own premium and its own risk measure. The premium is the price point of the cash
line, the sure payment at which the book changes hands, and the risk measure is loss ratio times premium,
visible even where the ratio explodes as the residue of the singularity. With this reading, a single
substantive axiom, comonotone pooling at relative cost, together with law invariance and regularity of the
implied pair, characterizes the loss ratios that are ratios of two Choquet-type functionals. The
characterization is sharp, and the pair is unique. At the level of representations it is a transfer of the
theory of \citet{yaari1987}, \citet{schmeidler1986}, \citet{wyp1997} and \citet{www2020}; at the level of the
loss ratio it is new, and it comes with facts that have no single-functional counterpart: no non-trivial loss
ratio is monotone (Proposition~\ref{prop:F2}), the association holds exactly off $\{\rho=\Pi=0\}$ and is
signed (Corollary~\ref{cor:assoc}), and pooling extends consistently to books with infinite or indeterminate
loss ratios (Corollary~\ref{cor:projpool}).

\paragraph{Open problems.}
\begin{enumerate}[label=(\arabic*)]
\item \emph{Beyond comonotone pooling.} The most promising direction is to drop \axCM{} in favour of an axiom
  that allows diversification, so that the implied premium and risk measure are convex or subadditive rather
  than comonotonic additive. The implied premium and the implied risk measure are defined without \axCM{} only
  partly (the premium needs only \axP; the accounting reading of $\rho$ needs pooling with sure losses), and the
  resulting class, a ratio of Kusuoka-type functionals, has no characterization yet. Replacing comonotonic
  additivity by convexity leads to suprema of signed Choquet integrals \citep{liu2020}, so no single pair $(h,g)$
  and no uniqueness should be expected.
\item \emph{Weaker cardinal axioms.} Can \axCM{} be replaced by the mediant inequality plus a supplementary
  condition? The inequality implies additivity for comonotone pairs strictly on the same side of break-even, but
  not for pairs straddling $1$ or lying in the break-even set.
\item \emph{Order properties of $\LR$.} \axRM{} is monotonicity of the implied pair. Whether it is equivalent to a
  natural order property of $\LR$ is open; by Proposition~\ref{prop:F2} it cannot be pointwise monotonicity.
\item \emph{Loadings on sure losses.} We conjecture that an analogue of Theorem~\ref{thm:A} holds for ratios
  $\bracket{I_h}{\kappa I_g}$ with a known loading $\kappa>0$, which would admit constant loss ratios $\ne1$.
\item \emph{Quasi-convexity.} Necessity of the conditions in Corollary~\ref{cor:shape}(d) is open.
\item \emph{Actuarial domain.} Whether (A4) imposed only on pairs with $X,Y,X+Y\notin N_0$ suffices is open; a
  truncation-continuity (Greco-type) version of the signed tier on the cone is not known.
\item \emph{Unbounded losses.} Extensions to $L^p$ and Orlicz spaces via \citet{www2020b} are not treated.
\end{enumerate}
